\documentclass[10pt]{article}
\usepackage[letterpaper,margin=0.7in]{geometry}
\usepackage{amsmath,amssymb}
\usepackage{booktabs}
\usepackage{parskip}
\setlength{\parskip}{4pt}
\pagestyle{empty}
\begin{document}

\begin{center}
{\large\bfseries Size distributions: what $N$, $N(D)$, $n(D)$ and $\mathrm{d}N/\mathrm{d}D$ mean, and what to fit}
\end{center}

\textbf{1. The quantities.} A probe counts particles in a sampled volume and sorts them into size bins
$i$ with lower and upper edges $D_i^{-}$, $D_i^{+}$, width $\Delta D_i = D_i^{+}-D_i^{-}$ (not constant
across bins) and a representative size $D_i$ (the bin midpoint).

\begin{center}
\begin{tabular}{@{}llll@{}}
\toprule
Symbol & Name & Units & Meaning \\
\midrule
$N_T$ & total number concentration & cm$^{-3}$ & all particles per unit volume of air \\
$\Delta N_i$ & bin concentration & cm$^{-3}$ & particles per unit volume in bin $i$ \\
$n(D) \equiv \dfrac{\mathrm{d}N}{\mathrm{d}D}$ & number distribution function & cm$^{-3}\,\mu$m$^{-1}$ &
particles per unit volume \emph{per unit diameter} \\
$f(D)$ & probability density of size & $\mu$m$^{-1}$ & $n(D)/N_T$; integrates to 1 \\
\bottomrule
\end{tabular}
\end{center}

The measured estimate of $n(D)$ in bin $i$ is the bin concentration divided by the bin width,
\begin{equation}
n(D_i) \approx \frac{\Delta N_i}{\Delta D_i}
\quad\Longleftrightarrow\quad
\Delta N_i \approx n(D_i)\,\Delta D_i ,
\end{equation}
which is why a code that stores \texttt{dNdD} recovers the per-bin concentration as
\texttt{dNdD * bin\_width}. The totals follow as a sum over bins or an integral over size:
\begin{equation}
N_T \;=\; \sum_i \Delta N_i \;=\; \int_0^\infty n(D)\,\mathrm{d}D .
\end{equation}
Integrating $\Delta N_i$ (cm$^{-3}$) over $D$ does \emph{not} give $N_T$; only $n(D)$ is integrated.

\textbf{2. The gamma form} (Johnson \& Kliche 2020, \emph{Atmosphere} 11, 333, Eqs.~1--2). The
distribution is written as total concentration times a probability density:
\begin{equation}
n(D) \;=\; N_T\, f(D;\mu,\lambda),
\qquad
f(D;\mu,\lambda) \;=\; \frac{\lambda^{\mu+1}}{\Gamma(\mu+1)}\, D^{\mu} e^{-\lambda D},\quad D>0 ,
\end{equation}
with $\int_0^\infty f\,\mathrm{d}D = 1$. Multiplying out gives the familiar three-parameter form
\begin{equation}
n(D) \;=\; N_0\, D^{\mu} e^{-\lambda D},
\qquad
N_0 \;=\; \frac{N_T\,\lambda^{\mu+1}}{\Gamma(\mu+1)},
\qquad
N_T \;=\; \int_0^\infty N_0 D^{\mu}e^{-\lambda D}\,\mathrm{d}D \;=\; \frac{N_0\,\Gamma(\mu+1)}{\lambda^{\mu+1}} .
\end{equation}
\begin{itemize}\setlength{\itemsep}{1pt}
\item The left side, often written $N(D)$, is the number distribution function $n(D)=\mathrm{d}N/\mathrm{d}D$
(cm$^{-3}\,\mu$m$^{-1}$), \emph{not} a per-bin concentration.
\item $N_0$ is the \emph{intercept} parameter, not the total concentration; its units,
cm$^{-3}\,\mu$m$^{-1-\mu}$, depend on $\mu$. It equals $N_T$ only in the special case $\lambda^{\mu+1}=\Gamma(\mu+1)$.
\item $\mu$ is the shape parameter (dimensionless); $\lambda$ is the \emph{rate} (slope) parameter
($\mu$m$^{-1}$), the inverse of the scale parameter $1/\lambda$.
\end{itemize}

\textbf{3. Size distribution versus probability distribution.} $n(D)$ is a concentration density: it
integrates to $N_T$, a physical number per volume. Dividing by $N_T$ gives $f(D)$, a probability density
of particle size that integrates to 1. So a gamma size distribution is a gamma probability density scaled
by $N_T$, exactly as Eq.~(3) writes it. By contrast $n(D)/N_0 = D^{\mu}e^{-\lambda D}$ integrates to
$\Gamma(\mu+1)/\lambda^{\mu+1}$, which is generally not 1.

\textbf{4. What to fit.} $x = D_i$ (bin midpoint) and $y = \mathrm{d}N/\mathrm{d}D$ at $D_i$, because the
model $N_0 D^{\mu}e^{-\lambda D}$ \emph{is} a density in $D$. Fitting the per-bin concentration
$\Delta N_i$ instead would make wide bins look more populated than narrow ones purely because they are
wide, and the shape of the fitted curve would follow the bin layout rather than the cloud. Equivalently,
one may fit $\Delta N_i$ against the model integrated over each bin,
$\int_{D_i^-}^{D_i^+} N_0 D^{\mu}e^{-\lambda D}\,\mathrm{d}D$, which is the more exact treatment for wide bins.

\vfill
{\footnotesize Notation in the literature is not uniform: some authors write $N(D)$ for what is called
$n(D)$ here. Read the units to tell which is meant: cm$^{-3}\,\mu$m$^{-1}$ is the distribution function,
cm$^{-3}$ is a bin concentration or a total.}

\end{document}
